
Large language models demonstrate substantial code generation capabilities, yet formal verification---the gold standard for correctness---remains challenging. When automated verifiers reject specifications for failing to prove safety or functional properties, current approaches typically discard the specification and regenerate from scratch, ignoring semantic signals encoded in verifier feedback that could guide systematic refinement.

This limitation creates a bottleneck in deploying formal verification to safety-critical systems. Formal specifications provide mathematical correctness guarantees essential for medical devices, aerospace control systems, and cryptographic implementations, but synthesizing specifications that verify requires expert verification engineers---a skillset rarer than software developers. Without automated specification synthesis, formal verification cannot scale to meet growing demand for provably correct software.

The challenge is not simply that LLMs struggle with formal reasoning. Rather, when verification fails, the semantic information in failure feedback is discarded rather than used to guide refinement. A failed proof obligation contains information: witness counterexamples showing where specifications fail, proof obligation structures revealing what needs proving, and dependency chains indicating why proofs fail. Yet this multi-dimensional semantic signal is either discarded entirely or presented to LLMs as unstructured natural language.

\textbf{Our Key Insight:} Verifier feedback can be viewed as a semantic gradient for specification synthesis. By decomposing feedback into three informational dimensions---Witness Instantiation, Logical Structure, and Dependency Preservation---we encode complementary semantic constraints that guide LLMs toward valid specifications via localized, targeted edits rather than global regeneration.

We demonstrate this insight through proof-of-concept experiments providing evidence for an information-theoretic framework for verification-in-loop. Our contributions are:

\begin{enumerate}
\item \textbf{Information-theoretic decomposition of verifier feedback:} We formalize three feedback dimensions and quantify their additive information value through empirical validation. Across four feedback conditions (RawError baseline, TagOnly, ObligationSlice, FullStructured), discharge rates scale monotonically from 31.9\% to 70.1\% with a linear information gradient ($\beta =12.49, R^{2}=0.89$, p<$10^{-50})$. Each dimension contributes independently, demonstrating non-redundant semantic constraints.
\end{enumerate}

\begin{enumerate}
\item \textbf{Cross-verifier semantic normalization via minimal taxonomy:} We introduce an 8-primitive taxonomy achieving 100\% error category coverage across Frama-C, Dafny, and Why3. This enables cross-verifier transfer with 15.1\% performance degradation, providing evidence that verifiers share a semantic core despite syntactic differences.
\end{enumerate}

\begin{enumerate}
\item \textbf{Causal evidence via compute-matched control:} Through controlled experiments isolating feedback quality from computational budget, we demonstrate that structured feedback drives systematic improvement beyond naive scaling. Under equal token budgets and verifier time, iterative feedback achieves 71.4\% discharge vs. 60.8\% for self-consistency sampling (10.7 percentage point gap, p<0.0001, Cohen's d=7.10).
\end{enumerate}

\begin{enumerate}
\item \textbf{Validation of non-vacuous specification strength:} Mutation testing demonstrates synthesized specifications achieve 63.3\% mutation kill rate, matching expert-written gold baseline (60\%), providing evidence of semantic meaningfulness beyond trivial ''spec washing.''
\end{enumerate}

\textbf{Scope and Limitations:} This proof-of-concept uses simulated verifier feedback (stochastic discharge rates 40-75\%) to control experimental variables and ensure reproducibility. Quantitative metrics represent proof-of-concept results requiring real-verifier validation. Our approach aligns with standard practices for mechanism validation.
